\documentclass[10pt,a4paper]{amsart}
\usepackage{amsmath,amssymb,amsthm,graphicx}
\usepackage{verbatim}
\newenvironment{reference}{\comment}{\endcomment}
%\newenvironment{reference}{}{}
\newenvironment{slogan}{\comment}{\endcomment}
\newenvironment{history}{\comment}{\endcomment}

% For commutative diagrams we use Xy-pic
\usepackage[all]{xy}

% We use 2cell for 2-commutative diagrams.
\xyoption{2cell}
% Theorem environments.
%
\theoremstyle{plain}
\newtheorem{theorem}[subsection]{Theorem}
\newtheorem{proposition}[subsection]{Proposition}
\newtheorem{lemma}[subsection]{Lemma}

\theoremstyle{definition}
\newtheorem{definition}[subsection]{Definition}
\newtheorem{example}[subsection]{Example}
\newtheorem{exercise}[subsection]{Exercise}
\newtheorem{situation}[subsection]{Situation}

\theoremstyle{remark}
\newtheorem{remark}[subsection]{Remark}
\newtheorem{remarks}[subsection]{Remarks}

\numberwithin{equation}{subsection}

% Macros
%
\def\lim{\mathop{\mathrm{lim}}\nolimits}
\def\colim{\mathop{\mathrm{colim}}\nolimits}
\def\Spec{\mathop{\mathrm{Spec}}}
\def\Hom{\mathop{\mathrm{Hom}}\nolimits}
\def\Ext{\mathop{\mathrm{Ext}}\nolimits}
\def\SheafHom{\mathop{\mathcal{H}\!\mathit{om}}\nolimits}
\def\SheafExt{\mathop{\mathcal{E}\!\mathit{xt}}\nolimits}
\def\Sch{\mathit{Sch}}
\def\Mor{\mathop{\mathrm{Mor}}\nolimits}
\def\Ob{\mathop{\mathrm{Ob}}\nolimits}
\def\Sh{\mathop{\mathit{Sh}}\nolimits}
\def\NL{\mathop{N\!L}\nolimits}
\def\CH{\mathop{\mathrm{CH}}\nolimits}
\def\proetale{{pro\text{-}\acute{e}tale}}
\def\etale{{\acute{e}tale}}
\def\QCoh{\mathit{QCoh}}
\def\Ker{\mathop{\mathrm{Ker}}}
\def\Im{\mathop{\mathrm{Im}}}
\def\Coker{\mathop{\mathrm{Coker}}}
\def\Coim{\mathop{\mathrm{Coim}}}

% Boxtimes
%
\DeclareMathSymbol{\boxtimes}{\mathbin}{AMSa}{"02}

%
% Macros for moduli stacks/spaces
%
\def\QCohstack{\mathcal{QC}\!\mathit{oh}}
\def\Cohstack{\mathcal{C}\!\mathit{oh}}
\def\Spacesstack{\mathcal{S}\!\mathit{paces}}
\def\Quotfunctor{\mathrm{Quot}}
\def\Hilbfunctor{\mathrm{Hilb}}
\def\Curvesstack{\mathcal{C}\!\mathit{urves}}
\def\Polarizedstack{\mathcal{P}\!\mathit{olarized}}
\def\Complexesstack{\mathcal{C}\!\mathit{omplexes}}
% \Pic is the operator that assigns to X its picard group, usage \Pic(X)
% \Picardstack_{X/B} denotes the Picard stack of X over B
% \Picardfunctor_{X/B} denotes the Picard functor of X over B
\def\Pic{\mathop{\mathrm{Pic}}\nolimits}
\def\Picardstack{\mathcal{P}\!\mathit{ic}}
\def\Picardfunctor{\mathrm{Pic}}
\def\Deformationcategory{\mathcal{D}\!\mathit{ef}}
\usepackage{cleveref}
\crefname{originallemma}{original lemma}{original lemmas}
\Crefname{originallemma}{Original lemma}{Original lemmas}
\makeatletter
\newlabel{algebra-lemma-separable-smooth}{{00TV}{source}{}{sourcepointer-1}{}}
\newlabel{algebra-lemma-separable-smooth@cref}{{[originallemma][0][]00TV}{[1][0][]source}{}{}{}}
\makeatother
\crefname{originalreference}{original reference}{original references}
\Crefname{originalreference}{Original reference}{Original references}
\makeatletter
\newlabel{lemma-nodal}{{0C4D}{source}{}{sourcepointer-2}{}}
\newlabel{lemma-nodal@cref}{{[originalreference][0][]0C4D}{[1][0][]source}{}{}{}}
\makeatother
\crefname{originallemma}{original lemma}{original lemmas}
\Crefname{originallemma}{Original lemma}{Original lemmas}
\makeatletter
\newlabel{more-algebra-lemma-completion-regular}{{07NY}{source}{}{sourcepointer-3}{}}
\newlabel{more-algebra-lemma-completion-regular@cref}{{[originallemma][0][]07NY}{[1][0][]source}{}{}{}}
\makeatother
\crefname{originallemma}{original lemma}{original lemmas}
\Crefname{originallemma}{Original lemma}{Original lemmas}
\makeatletter
\newlabel{properties-lemma-number-of-branches-irreducible-components}{{0E20}{source}{}{sourcepointer-4}{}}
\newlabel{properties-lemma-number-of-branches-irreducible-components@cref}{{[originallemma][0][]0E20}{[1][0][]source}{}{}{}}
\makeatother
\crefname{originalreference}{original reference}{original references}
\Crefname{originalreference}{Original reference}{Original references}
\makeatletter
\newlabel{remark-trivial-quadratic-extension}{{0CBU}{source}{}{sourcepointer-5}{}}
\newlabel{remark-trivial-quadratic-extension@cref}{{[originalreference][0][]0CBU}{[1][0][]source}{}{}{}}
\makeatother
\title{Reduced closed subschemes of a nodal curve are nodal, and meet their complement in finite separable spectra at nodes of X}
\author{The Stacks Project Authors}
\date{stacks-project:curves.tex:local-snapshot-20261012; GNU FDL 1.2 or later}
\begin{document}
\maketitle
\noindent\textit{Editorial note.} All delivered mathematical text is the author's own, copied from the pinned Stacks Project chapter file curves.tex; no statement or proof step was rewritten, reordered or removed.\par\medskip
\noindent\textit{Editorial note.} The counted proof is the author's complete argument: the component-union reduction, the full case analysis including the two-branch situation, the completed uv/(uv) computation of the intersection ideal, and the closing regularity-plus-separability smoothness argument are all retained verbatim.\par\medskip
\noindent\textit{Editorial note.} The label lemma-closed-subscheme-nodal-curve names the counted statement, and the proof environment immediately following it in the source is delivered verbatim as its proof.\par\medskip
\noindent\textit{Editorial note.} The counted statement is carried by the author's own lemma environment in the chapter source.\par\medskip
\noindent\textit{Allowed standard prerequisites (editorial).}\begin{itemize}
\item Components through a node. Properties, Lemma properties-lemma-number-of-branches-irreducible-components; Same chapter, Lemma lemma-nodal. Bounds the number of irreducible components through a point by the number of geometric branches, which is 2 at a node.
\item The completed uv/(uv) normal form at a node. Same chapter, Remark remark-trivial-quadratic-extension and Lemma lemma-nodal; More on Algebra, Lemma more-algebra-lemma-completion-regular; Algebra, Lemma algebra-lemma-separable-smooth. Identifies the completed local ring as kappa(y)[[u, v]]/(uv), makes the completed maps to Y and Z the two coordinate projections, and upgrades regularity plus separability to smoothness.
\end{itemize}\medskip
\ifdefined\hypertarget\hypertarget{sourcepointer-1}{}\fi
\noindent\textit{Source reference.} Original lemma 00TV: sources/stacks-chow/tags.txt, source lines 1089--1089. Prerequisite pointer invoked at proof line(s) 4858; it names a standard fact of a sibling chapter that the counted argument uses, and every step of the counted proof built on it is retained verbatim in the delivered body.\par\medskip
\ifdefined\hypertarget\hypertarget{sourcepointer-2}{}\fi
\noindent\textit{Source reference.} Original reference 0C4D: sources/stacks-curves/pinned-original-source.tex, source lines 4391--4391. Prerequisite pointer invoked at proof line(s) 4826; it names a standard fact of the same chapter that the counted argument uses, and every step of the counted proof built on it is retained verbatim in the delivered body.\par\medskip
\ifdefined\hypertarget\hypertarget{sourcepointer-3}{}\fi
\noindent\textit{Source reference.} Original lemma 07NY: sources/stacks-chow/tags.txt, source lines 9778--9778. Prerequisite pointer invoked at proof line(s) 4855; it names a standard fact of a sibling chapter that the counted argument uses, and every step of the counted proof built on it is retained verbatim in the delivered body.\par\medskip
\ifdefined\hypertarget\hypertarget{sourcepointer-4}{}\fi
\noindent\textit{Source reference.} Original lemma 0E20: sources/stacks-chow/tags.txt, source lines 17590--17590. Prerequisite pointer invoked at proof line(s) 4829; it names a standard fact of a sibling chapter that the counted argument uses, and every step of the counted proof built on it is retained verbatim in the delivered body.\par\medskip
\ifdefined\hypertarget\hypertarget{sourcepointer-5}{}\fi
\noindent\textit{Source reference.} Original reference 0CBU: sources/stacks-curves/pinned-original-source.tex, source lines 4549--4549. Prerequisite pointer invoked at proof line(s) 4840; it names a standard fact of the same chapter that the counted argument uses, and every step of the counted proof built on it is retained verbatim in the delivered body.\par\medskip
\begin{lemma}
\label{lemma-closed-subscheme-nodal-curve}
Let $k$ be a field. Let $X$ be a locally algebraic $k$-scheme
equidimensional of dimension $1$ whose singularities are at-worst-nodal.
If $Y \subset X$ is a reduced closed subscheme
equidimensional of dimension $1$, then
\begin{enumerate}
\item the singularities of $Y$ are at-worst-nodal, and
\item if $Z \subset X$ is the scheme theoretic closure of
$X \setminus Y$, then
\begin{enumerate}
\item the scheme theoretic intersection $Y \cap Z$ is
the disjoint union of spectra of finite separable extensions of $k$,
\item each point of $Y \cap Z$ is a node of $X$, and
\item $Y \to \Spec(k)$ is smooth at every point of $Y \cap Z$.
\end{enumerate}
\end{enumerate}
\end{lemma}

\begin{proof}
Since $X$ and $Y$ are reduced and equidimensional of dimension $1$,
we see that $Y$ is the scheme theoretic union of a subset of the
irreducible components of $X$ (in a reduced ring $(0)$
is the intersection of the minimal primes).
Let $y \in Y$ be a closed point.
If $y$ is in the smooth locus of
$X \to \Spec(k)$, then $y$ is on a unique irreducible component
of $X$ and we see that $Y$ and $X$ agree in an open neighbourhood
of $y$. Hence $Y \to \Spec(k)$ is smooth at $y$.
If $y$ is a node of $X$ but still lies on a unique irreducible
component of $X$, then $y$ is a node on $Y$ by the same argument.
Suppose that $y$ lies on more than $1$ irreducible component of $X$.
Since the number of geometric branches of $X$ at $y$ is $2$
by Lemma \ref{lemma-nodal},
there can be at most $2$ irreducible components passing through $y$
by Properties, Lemma
\ref{properties-lemma-number-of-branches-irreducible-components}.
If $Y$ contains both of these, then again $Y = X$ in an open neighbourhood
of $y$ and $y$ is a node of $Y$. Finally, assume $Y$ contains only one
of the irreducible components. After replacing $X$ by an open
neighbourhood of $x$ we may assume $Y$ is one of the two
irreducble components and $Z$ is the other. By Properties, Lemma
\ref{properties-lemma-number-of-branches-irreducible-components}
again we see that $X$ has two branches at $y$, i.e., the local ring
$\mathcal{O}_{X, y}$ has two branches and that these branches
come from $\mathcal{O}_{Y, y}$ and $\mathcal{O}_{Z, y}$. Write
$\mathcal{O}_{X, y}^\wedge \cong \kappa(y)[[u, v]]/(uv)$
as in Remark \ref{remark-trivial-quadratic-extension}.
The field $\kappa(y)$ is finite separable over $k$ by
Lemma \ref{lemma-nodal} for example.
Thus, after possibly switching the roles of $u$ and $v$,
the completion of the map
$\mathcal{O}_{X, y} \to \mathcal{O}_{Y, Y}$
corresponds to $\kappa(y)[[u, v]]/(uv) \to \kappa(y)[[u]]$
and the completion of the map
$\mathcal{O}_{X, y} \to \mathcal{O}_{Y, Y}$
corresponds to $\kappa(y)[[u, v]]/(uv) \to \kappa(y)[[v]]$.
The scheme theoretic intersection of $Y \cap Z$ is cut out
by the sum of their ideas which in the completion is $(u, v)$, i.e.,
the maximal ideal. Thus (2)(a) and (2)(b) are clear.
Finally, (2)(c) holds: the completion of $\mathcal{O}_{Y, y}$
is regular, hence $\mathcal{O}_{Y, y}$ is regular
(More on Algebra, Lemma \ref{more-algebra-lemma-completion-regular})
and $\kappa(y)/k$ is separable, hence
smoothness in an open neighbourhood
by Algebra, Lemma \ref{algebra-lemma-separable-smooth}.
\end{proof}

\end{document}
